\documentclass[a4paper,11pt,lualatex,ja=standard]{bxjsarticle}
\geometry{top=25mm,bottom=25mm,left=25mm,right=25mm}
\usepackage{graphicx}
\usepackage{booktabs}
\usepackage{array}
\usepackage{longtable}
\usepackage{xcolor}

% 研究の基本情報を読み込み（protocol_template.texと共通）
\input{research_info.tex}

\begin{document}

\begin{center}
{\LARGE \textbf{【\ResearchTitle】}}

\vspace{1em}

{\Large \textbf{説明同意文書}}

\vspace{2em}

説明日：\hspace{3em}年\hspace{2em}月\hspace{2em}日

\vspace{2em}

{\small \ProtocolVersion\hspace{1em}\ProtocolDate}
\end{center}

\vspace{2em}

\section{研究実施について}

「\ResearchTitle」は\InstitutionFullName\ 医の倫理委員会の審査を受け、研究機関の長の許可を受けて実施しています。

\section{研究機関}

\begin{itemize}
\item 名称：\InstitutionFullName
\item 研究責任者氏名：%%% PLACEHOLDER:PI_NAME %%%
\end{itemize}

% 共同研究機関がある場合は以下のコメントを外してください
% \begin{itemize}
% \item 共同研究機関 名称：\CollaboratingInstitution
% \item 共同研究機関 研究責任者氏名：\CollaboratingPI
% \end{itemize}

\section{研究の目的および意義}

ここに研究の目的とその研究成果が医学・医療にもたらす価値を簡潔に記載してください。

\begin{itemize}
\item 研究の背景と現在何がわかっているか
\item 何を明らかにしようとしているか
\item 社会的および学術的な意義
\end{itemize}

\textbf{【記載例】}

本研究の目的は、●●病の病態解明と新しい治療法の開発に寄与することです。現在、●●病については...ということがわかっていますが、...については明らかになっていません。このことが明らかになることで、...に役立つと考えられます。

\section{研究方法と期間}

\subsection{研究方法}

\begin{itemize}
\item 研究対象者から取得する試料・情報の利用目的および取り扱い
\item いつ、どのようにして試料を採取するか（カルテ調査や質問紙調査を含む）
\item 調べる内容
\item 研究のスケジュール
\end{itemize}

\textbf{【記載例】}

この研究では、以下の方法で実施します：

\begin{enumerate}
\item 血液検査：XX mlの採血を行います
\item アンケート調査：XX項目についてお答えいただきます
\item カルテ調査：過去の診療記録を参照させていただきます
\end{enumerate}

\subsection{研究期間}

\begin{itemize}
\item 研究実施期間：倫理審査承認日から令和XX年XX月XX日まで
\item 研究対象者登録期間：令和XX年XX月XX日から令和XX年XX月XX日まで
\end{itemize}

\section{研究対象者として選定された理由}

\begin{itemize}
\item 対象者に協力いただく必要性
\item 適格基準・除外基準
\item 対象者の人数
\end{itemize}

\textbf{【記載例】}

本研究では●●病と診断された患者さんXX名を対象としています。あなたは●●の理由により、本研究の対象者として選定されました。

\section{研究対象者に生じる負担と、予想されるリスクおよび利益}

\subsection{負担とリスク}

\begin{itemize}
\item 採血や検査などの身体的負担
\item 時間的な負担
\item 予測される身体的・心理的・社会的リスク
\end{itemize}

\textbf{【記載例】}

\begin{itemize}
\item 採血による痛みや内出血の可能性があります
\item アンケートへの回答に約XX分の時間を要します
\item ...
\end{itemize}

\subsection{利益}

\begin{itemize}
\item 研究参加により得られる可能性のある利益
\item 将来の患者さんへの貢献
\end{itemize}

\textbf{【記載例】}

本研究に参加することで、あなた自身に直接的な利益があるわけではありませんが、この研究によって●●病の治療法の開発が進み、将来の患者さんたちに役立つ可能性があります。

\section{いつでも同意の撤回ができます}

参加を決めた後も、いつでも参加をやめることができます。

\textbf{【記載例】}

参加を決めて採血を行った後も、いつでも参加をやめることができます。参加をやめる場合は、血液やそれまでの調査記録は破棄します。ただし、お申し出があった時にすでに研究結果が公表されていたときなど、データから除けない場合もあります。

参加をやめる場合には、文書を書いていただきますので、下記の研究者までご連絡ください。

\section{研究に同意しない、または同意撤回において不利益はありません}

研究へ参加するかどうか、もしくは継続するかどうかはよくお考えのうえ、自由に決めてください。同意しない、もしくは同意を撤回された場合も、不利益な扱いを受けることは一切なく、そのときの最善の治療を行います。

\section{研究に関する情報公開について}

この研究は学術雑誌や学会での発表を予定しています。研究の進行状況は...でお知らせいたします。

\section{研究計画書等の閲覧について}

研究について詳しく知りたい場合は、他の研究対象者の個人情報保護や研究の独創性に支障のない範囲で研究計画書や研究の方法に関する資料を見ることができます。下記の研究者までお問い合わせください。

\section{個人情報の取扱いについて}

対象者の方々の個人情報の保護には十分配慮いたします。

\begin{itemize}
\item 個人情報は暗号化され、結果をお返しするとき以外は番号で扱われます
\item 番号の対応表や同意書などの研究に関わる書類やデータ、試料は厳重に保管します
\item 研究の結果は学術雑誌や学会発表で公表する予定ですが、この時にも個人の情報が使用されることはありません
\end{itemize}

\textbf{匿名化の方法：}

個人を特定できる情報（氏名、住所など）は研究ID番号に置き換えられ、厳重に管理されます。

\section{試料・情報の保管および廃棄の方法}

\subsection{保管方法}

研究で得られた試料および情報は施錠された書庫や冷蔵庫で厳重に保管します。

\subsection{保管期間}

研究終了後10年間、試料・情報は保管し、その後、廃棄します。

ただし、後に説明する「同意をうける時点では想定されない将来の研究」に使用するため、長期間にわたり保管する可能性があります。

\subsection{廃棄方法}

保管期間終了後は、個人情報に配慮し、適切に廃棄します。

\section{研究資金および利益相反について}

\textbf{【記載例1：運営費交付金で実施する場合】}

本研究は、運営費交付金により実施します。また、本研究は、特定の企業からの資金提供を受けていません。本研究の実施にあたり、利益相反については、「京都大学利益相反ポリシー」「京都大学利益相反マネジメント規程」に従い、「京都大学臨床研究利益相反審査委員会」において適切に審査しています。

\textbf{【記載例2：公的研究費で実施する場合】}

本研究は、公的研究費である○○省科学研究費（○○○○）により実施します。また、本研究は、特定の企業からの資金提供を受けていません。本研究の実施にあたり、利益相反については、「京都大学利益相反ポリシー」「京都大学利益相反マネジメント規程」に従い、「京都大学臨床研究利益相反審査委員会」において適切に審査しています。

\section{研究より得られた結果の取り扱い}

\subsection{結果の開示について}

検査結果がわかるまでには、...週間ほどかかる予定です。結果はご本人にのみにお伝えし、たとえご家族であってもご本人の承諾なしには説明することができませんのでご了解ください。

また、結果開示をしない選択や、希望した時期に知らせる方法もあります。途中で気持ちが変わったときは結果を聞かないこともできます。

\textbf{注意：}これは研究の結果であり、診療の用に供する検査ではありません。このため、診療の用に供する場合に求められる精度管理が制度的に行われているものではないことに留意の上、適切にお取扱いください。

\section{相談の窓口}

説明の中でわからない言葉や質問、また参加や結果開示のことで相談がありましたら何でも遠慮せずにお話しください。

\subsection{研究課題ごとの相談窓口}

%%% PLACEHOLDER:CONSULTATION_CONTACT %%%

\subsection{京都大学の苦情等の相談窓口}

\begin{enumerate}
\item 研究対象者が京大病院の患者の場合、京大病院の教職員が行う研究の場合：
  \begin{itemize}
  \item 京都大学医学部附属病院 臨床研究相談窓口
  \item Tel: 075-751-4748
  \item E-mail: ctsodan@kuhp.kyoto-u.ac.jp
  \end{itemize}

\item 京大病院が関与しない医学研究科の研究の場合：
  \begin{itemize}
  \item 京都大学医学研究科 総務企画課 研究推進掛
  \item Tel: 075-753-9301
  \item E-mail: 060kensui@mail2.adm.kyoto-u.ac.jp
  \end{itemize}

\item 遺伝カウンセリングに関する窓口（必要に応じて）：
  \begin{itemize}
  \item 京都大学医学部附属病院 遺伝子診療部
  \item Tel: 075-751-4350（受付時間 平日 13:00〜16:30）
  \end{itemize}
\end{enumerate}

\section{経済的負担／謝礼について}

\textbf{【記載例1：負担がない場合】}

本研究は研究費で実施するため、対象者のご負担はありません。ただし、外来受診については通常の保険診療内でご負担いただきます。

\textbf{【記載例2：謝礼がある場合】}

本研究では、対象者おひとりにつきXX円の謝礼をお支払いします。別途、必要に応じて交通費をお支払いいたしますので、所定の用紙にて手続きをお願いいたします。

\section{試料・情報の将来の研究における使用および他機関への提供}

本研究で収集した試料・情報は、同意を受ける時点では特定されない将来の研究のために用いる可能性があります。他の研究への二次利用および他研究機関へ提供する際は、新たな研究計画について倫理審査委員会で承認された後に行います。また、ホームページ上で、研究の目的を含む研究実施の情報を公開し、研究対象者が拒否できる機会を保障します。

% 以下の項目は該当する場合のみ含めてください

% \section{（通常の診療を超える医療行為を伴う研究の場合）他の治療法について}
%
% 本研究に参加しない場合も、現状における最善の治療法を実施します。具体的には...

% \section{（通常の診療を超える医療行為を伴う研究の場合）研究参加後の医療の提供について}
%
% 研究実施後には XX や XX といった診療を実施する可能性があります。

% \section{（侵襲を伴う研究の場合）健康被害の補償について}
%
% 本研究参加に伴い、健康被害が発生した場合には、通常診療にて対応いたします。診療は保険診療の範囲内で実施し、費用負担が発生することになります。

% \section{（侵襲を伴う研究であって介入を行うものの場合）モニタリング・監査・倫理審査委員会による個人情報の閲覧}
%
% 本研究においては、研究の進捗状況や計画通り実施されているかを確認する「モニタリング」と、適正に研究が実施されたかを確認する「監査」を受けます。これらの業務に必要な範囲で、モニタリングや監査に従事する者、および倫理審査委員会が、情報や秘密の保全を前提として、対象者に関する試料や情報を閲覧することがあります。

\vfill

この研究についてご理解していただき、参加していただける場合は「研究参加の同意書」に署名していただきます。この説明文書は差し上げますので、よく読んでご検討ください。

\vspace{1em}

{\small \textbf{※省略した項目がある場合、項目の通し番号に注意してください。}}

\end{document}
